% Supplementary material for main_grsl_revised.tex. Compiled separately.
% Build: latexmk -pdf -interaction=nonstopmode -halt-on-error supplementary.tex
% Tables are generated by experiments/figures/make_supp_tables.py.
\documentclass[journal]{IEEEtran}

\usepackage{graphicx}
\usepackage{booktabs}
\usepackage{multirow}
\usepackage{amsmath}
\usepackage{amssymb}
\usepackage{cite}
\usepackage[hidelinks]{hyperref}
\makeatletter\newcommand{\tablebody}[1]{\@@input #1 }\makeatother

\begin{document}

\title{Supplementary Material\\Boundary Veto: Model-Free Post-Processing for Agricultural Field Boundary Delineation}
\author{Aviskar~Giri and Isaac~Corley}
\maketitle

All results use the letter's prediction runs and released reference polygons.

\section{Full Parameter Sweeps}
PQ remains flat for support thresholds between 0.7 and 0.9 on every model and sensor pair, while $h=8$ lowers accuracy because small fields lose their seeds.
We report every setting of $\tau$ and $h$ and each pair's foreground and split-rule comparison in Table~\ref{tab:supp_sweep}, using the letter's macro average.
The letter includes the PlanetScope B3 and Sentinel-2 B7 columns.

\begin{table}[h]
\centering
\caption{Accuracy Holds Near the Chosen Thresholds}
\label{tab:supp_sweep}
\footnotesize
\setlength{\tabcolsep}{4pt}
\begin{tabular}{@{}lcccc@{}}
\toprule
Setting & PS B3 & PS B7 & S2 B7 & S2 B3 \\
\midrule
\tablebody{tables/supp_sweep_rows.tex}
\bottomrule
\end{tabular}
\end{table}

\section{Complete Grids for Prior Methods}
We choose each prior method's parameters by validation PQ on PlanetScope B3, then apply the 900~m$^2$ minimum-size step for the letter's comparison.
Table~\ref{tab:supp_priors} includes every setting we tried and its validation and test scores.
The best settings for probability thresholds, iterative marker depth, and stochastic watershed lie at the search grid's edge.
Iterative marker depth and stochastic watershed stay more than nine points below the FTP conversion across all settings tried.

\begin{table}[h]
\centering
\caption{No Setting of Any Prior Conversion Reaches FTP}
\label{tab:supp_priors}
\footnotesize
\setlength{\tabcolsep}{5pt}
\begin{tabular}{@{}lcc@{}}
\toprule
Setting & Validation PQ & Test PQ \\
\midrule
\tablebody{tables/supp_priors_rows.tex}
\bottomrule
\end{tabular}
\smallskip

\parbox{\columnwidth}{\footnotesize
\texttt{thr-te$x$-tb$y$}: field probability at least $x$, boundary probability below $y$.
\texttt{ucm-k$x$}: merge borders whose mean boundary probability is below $x$.
\texttt{iterh-a$x$}: re-segment regions over $x$ pixels, on the field (\texttt{arg}) or not-background (\texttt{nbg}) foreground.
\texttt{swst-n$x$-m$y$-t$z$}: $y$ draws of $x$ markers, cut where the line frequency reaches $z$.}
\end{table}

\section{Alternatives Named in the Ablation}
We tested the foreground alternatives without splitting, replacing the assignment to the most probable class.
For the added steps, we used the boundary veto at $\tau=0.7$ and $h=2$ without the minimum-size step and scored only the validation split.
Table~\ref{tab:supp_variants} reports PlanetScope B3 validation PQ, and test PQ where scored, for every alternative mentioned without a number in the letter's ablation.
Pixel counts refer to the 3~m PlanetScope grid.

\begin{table}[h]
\centering
\caption{Every Alternative to the Veto Scores Lower}
\label{tab:supp_variants}
\footnotesize
\setlength{\tabcolsep}{3pt}
\begin{tabular}{@{}lcc@{}}
\toprule
Variant & Validation PQ & Test PQ \\
\midrule
\tablebody{tables/supp_variants_rows.tex}
\bottomrule
\end{tabular}
\end{table}

\section{Scores by Field Size}
The boundary veto gains most on small fields in every run, between 3.5 and 4.8 PQ points, and least on large fields, between 1.1 and 1.7 (Table~\ref{tab:supp_size}).
We group PQ and RQ by reference field area using the Fields of the Planet evaluation code's bins, with small fields below 0.5~ha, medium fields from 0.5 to 2~ha, and large fields above 2~ha.
On PlanetScope B3, the ten countries contain 29,406 small, 20,884 medium, and 12,173 large reference fields, making the small bin the largest.

\begin{table}[h]
\centering
\caption{The Gain Concentrates in the Smallest Fields}
\label{tab:supp_size}
\footnotesize
\setlength{\tabcolsep}{3.4pt}
\begin{tabular}{@{}lcccccc@{}}
\toprule
 & \multicolumn{3}{c}{PQ} & \multicolumn{3}{c}{RQ} \\
\cmidrule(lr){2-4}\cmidrule(lr){5-7}
Method & $<$0.5 ha & 0.5--2 & $>$2 & $<$0.5 ha & 0.5--2 & $>$2 \\
\midrule
\tablebody{tables/supp_size_rows.tex}
\bottomrule
\end{tabular}
\end{table}

\section{Per-Country Scores}
We compare the FTP conversion and the boundary veto for each country and model in Table~\ref{tab:supp_country}, giving the scores plotted in Fig.~3 of the letter.
Portugal has two test patches and Cambodia sixteen, so their scores use far fewer fields than the others.

\begin{table}[h]
\centering
\caption{The Veto Leads in Almost Every Country and Model}
\label{tab:supp_country}
\footnotesize
\setlength{\tabcolsep}{3pt}
\begin{tabular}{@{}lcccccccc@{}}
\toprule
 & \multicolumn{2}{c}{PS B3} & \multicolumn{2}{c}{PS B7} & \multicolumn{2}{c}{S2 B7} & \multicolumn{2}{c}{S2 B3} \\
\cmidrule(lr){2-3}\cmidrule(lr){4-5}\cmidrule(lr){6-7}\cmidrule(lr){8-9}
Country & FTP & Veto & FTP & Veto & FTP & Veto & FTP & Veto \\
\midrule
\tablebody{tables/supp_country_rows.tex}
\bottomrule
\end{tabular}
\end{table}

\section{Timing by Parcel Density}
We sampled thirty patches per bin with a fixed random seed and timed them on one core, restricting all threads to that core.
We kept the fastest of three repeats after a warm-up call and report median milliseconds per patch in four field-count bins on both sensors (Table~\ref{tab:supp_timing}).
Sentinel-2 patches contain more pixels, so every method takes longer there.
The ordering between methods is unchanged in the sparsest and densest bins, and iterative marker depth is slower than the boundary veto in the two middle bins.

\begin{table}[h]
\centering
\caption{The Veto Costs About What the Watershed Costs}
\label{tab:supp_timing}
\footnotesize
\setlength{\tabcolsep}{2.3pt}
\begin{tabular}{@{}lcccccccc@{}}
\toprule
 & \multicolumn{4}{c}{PS B3} & \multicolumn{4}{c}{S2 B7} \\
\cmidrule(lr){2-5}\cmidrule(lr){6-9}
Method & 1--9 & 10--29 & 30--99 & 100+ & 1--9 & 10--29 & 30--99 & 100+ \\
\midrule
\tablebody{tables/supp_timing_rows.tex}
\bottomrule
\end{tabular}
\end{table}

\section{Reproduction}
We save each released model's class probabilities once with its released inference settings, then use those same predictions for every conversion.
To reproduce the comparison, convert the saved probabilities to polygons with each method, score them against the reference polygons, and generate the tables and figures.
Our FTP conversion implementation reproduces published per-country PQ to within $10^{-3}$ for PlanetScope and for Sentinel-2 on the 10~m grid used by that paper's evaluation.

\end{document}
